\documentclass[12pt,a4paper,twoside]{book}
\usepackage{graphicx}
\usepackage{setspace}	%double spacing for text, single for captions, footnotes, etc.
%\usepackage{hypernat} 	%substitut de cite que permet fer hyperlinks
\usepackage{natbib}		% substituye a 'hypernat' que funciona en Windows.
\usepackage[spanish]{babel}
\usepackage[utf8]{inputenc}
\usepackage{color}
\usepackage{hhline} 		% extended styles for tables
\usepackage{multirow}
\usepackage{subfigure}
\usepackage{acronym}
\usepackage{hyperref}
\usepackage{amsmath,amsmath,amssymb} 
\usepackage{fancyhdr}
\usepackage{epsfig, amsmath}
\usepackage{algorithm}
\usepackage{algorithmic}

% general settings
\hypersetup{
	linktocpage=true,
	colorlinks=true,
	linkcolor=blue,
	citecolor=blue,
}
\definecolor{Hgray}{gray}{0.6}

\newenvironment{definition}[1][Definition]{\begin{trivlist}
\item[\hskip \labelsep {\bfseries #1}]}{\end{trivlist}}

\setlength{\topmargin}{0cm}
\setlength{\textheight}{23cm}
\setlength{\textwidth}{17cm}
\setlength{\oddsidemargin}{0cm}
\setlength{\evensidemargin}{0cm}
\setlength{\headheight}{1cm}

% indica que las 'sub-sub-sections' sean numeradas y aparezcan en el indice
\setcounter{secnumdepth}{3}
\setcounter{tocdepth}{2}

% settings for code
\renewcommand{\algorithmicrequire}{\textbf{Entrada: }}
\renewcommand{\algorithmicensure}{\textbf{Salida: }}

%%%%%%%%%%%%
% DOCUMENT %
%%%%%%%%%%%%
\begin{document}

% portada
\newpage
\thispagestyle{empty}

\baselineskip 2em

%\vspace*{1cm}

\centerline{\includegraphics[width=0.6\textwidth]{images/UOC-logo}}
\begin{center}
\textsc{Universitat Oberta de Catalunya (UOC) \\
 Máster Universitario en Ciencia de Datos (\textit{Data Science})\\}

%\centerline {\pic{UOC}{4cm}}

\vspace*{1.5cm}

\textsc{\Large TRABAJO FINAL DE MÁSTER}

\vspace*{0.5cm}

\textsc{\large Área: 2}


%\textbf{\Huge VirtualTechLab Model: }

\vspace*{2.0cm}

\textbf{\Large CLASIFICADOR DOCUMENTOS MÉDICOS HOPE }

\vspace{2.5cm}
\baselineskip 1em

\baselineskip 2em
-----------------------------------------------------------------------------\\
Autor:      Rubén Vasallo Gonzalez \\
Profesor:   Jordi Casas Roma\\
Tutor:      Carlos Luis Sanchez Bocanegra\\
Co-Tutor:   Rafael Pastor Vargas\\
-----------------------------------------------------------------------------\\
\vspace*{1.5cm}
Barcelona, \today

\end{center}

\newpage
\pagestyle{empty}
\hfill

\newpage
% abstract
\pagenumbering{roman} 
\setcounter{page}{1} 
\pagestyle{plain}

%%%%%%%%%%%%%%%%
%%% CREDITOS %%%
%%%%%%%%%%%%%%%%
\chapter*{Créditos/Copyright}

\vspace{1cm}

\begin{figure}[ht]
    \centering
	\includegraphics[scale=1]{images/license.png}
\end{figure}

Esta obra está sujeta a una licencia de Reconocimiento -  NoComercial - SinObraDerivada

\href{https://creativecommons.org/licenses/by-nc-nd/3.0/es/}{3.0 España de CreativeCommons}.

%%%%%%%%%%%%%
%%% FICHA %%%
%%%%%%%%%%%%%
\chapter*{FICHA DEL TRABAJO FINAL}

\begin{table}[ht]
	\centering{}
	\renewcommand{\arraystretch}{2}
	\begin{tabular}{r | l}
		\hline
		Título del trabajo: & CLASIFICADOR DOCUMENTOS MÉDICOS HOPE\\
		\hline
        Nombre del autor: & Rubén Vasallo Gonzalez\\
		\hline
        Nombre del colaborador/a docente: & Carlos Luis Sanchez Bocanegra, Rafael Pastor Vargas \\
		\hline
        Nombre del PRA: & Jordi Casas Roma \\
		\hline
        Fecha de entrega (mm/aaaa): & 01/2021\\
		\hline
        Titulación o programa: & Máster Universitario en Ciencia de Datos\\
		\hline
        Área del Trabajo Final: & M2.979 - TFM\\
		\hline
        Idioma del trabajo: & Español\\
		\hline
        Palabras clave & hope, clasificador, medicina\\
		\hline
	\end{tabular}
\end{table}

%%%%%%%%%%%%%%%%%%%
%%% DEDICATORIA %%%
%%%%%%%%%%%%%%%%%%%
\chapter*{Dedicatoria/Cita}

Quiero dedicarle este trabajo a mis mentores \textit{Carlos Luis Sanchez Bocanegra} y \textit{Rafael Pastor Vargas} que me han ayudado, apoyado y guiado en todo momento para conseguir los objetivos del máster.

%%%%%%%%%%%%%%%%%%%
%%% Agradecimientos %%%
%%%%%%%%%%%%%%%%%%%
\chapter*{Agradecimientos}

\paragraph{}
Quiero agradecer a \textit{Carlos Luis Sanchez Bocanegra} por invitarme participar en el Proyecto HOPE y poder aportar mi granito de arena a este gran proyecto.

\paragraph{}
También quiero dar la gracias a todos los miembros del proyecto HOPE que me han dado la bienvenida al grupo y me han facilitado la vida en unas circunstancias en las que, cuando entre en este, no eran las más idóneas. La mayoría de integrantes del grupo son médicos y la saturación de trabajo que había por el COVID-19 era enorme. Tengo claro que sin ellos y sin la ayuda en especial de Carlos y \textit{Nicolas Passadore} que ha estado luchando para conseguir un dataset con más observaciones, este Máster no habría sido posible.

%%%%%%%%%%%%%%%%
%%% RESUMEN  %%%
%%%%%%%%%%%%%%%%
\chapter*{Abstract}
\addcontentsline{toc}{chapter}{Abstract}

\onehalfspacing
\paragraph{}
This Final Master's Thesis {\textit{FMT}} was born from the need to be able to have in a simple, up-to-date and immediate way, medical bibliographic references cataloged according to the information and the patient's symptoms, being able to make a ranking of more or less interest function of the feedback provided by health professionals on these bibliographic references.

\paragraph{}
\textbf{Resumen}:
\paragraph{}
Este Trabajo de Final de Máster (\textit{TFM}) nace de la necesidad de poder disponer de una manera sencilla, actualizada e inmediata, referencias bibliográficas médicas catalogadas según la información y los síntomas que tienen del paciente, pudiendo hacer una clasificación (\textit{ranking}) de más o menos interés en función de la valoración (\textit{feedback}) aportada por los profesionales sanitarios sobre estas referencias bibliográficas.

\vspace{1.5cm}

\textbf{Palabras clave}: hope, clasificador, artículos, médicos, PCA, KNN, Regresión Logística, logit, Random Forest, SVM
\newpage

\pagestyle{fancy}
\renewcommand{\chaptermark}[1]{ \markboth{#1}{}}
\renewcommand{\sectionmark}[1]{\markright{ \thesection.\ #1}}
\lhead[\fancyplain{}{\bfseries\thepage}]{\fancyplain{}{\bfseries\rightmark}}
\rhead[\fancyplain{}{\bfseries\leftmark}]{\fancyplain{}{\bfseries\thepage}}
\cfoot{}

% indice
\cleardoublepage
\phantomsection
\addcontentsline{toc}{chapter}{Índice}
\tableofcontents
% listado de figuras
\cleardoublepage
\phantomsection
\addcontentsline{toc}{chapter}{Llistado de Figuras}
\listoffigures
% listado de tablas
\cleardoublepage
\phantomsection
\addcontentsline{toc}{chapter}{Listado de Tablas}
\listoftables

\thispagestyle{empty}

\pagenumbering{arabic}

\pagestyle{fancy}
\renewcommand{\chaptermark}[1]{ \markboth{#1}{}}
\renewcommand{\sectionmark}[1]{\markright{ \thesection.\ #1}}
\lhead[\fancyplain{}{\bfseries\thepage}]{\fancyplain{}{\bfseries\rightmark}}
\rhead[\fancyplain{}{\bfseries\leftmark}]{\fancyplain{}{\bfseries\thepage}}
\cfoot{}

\onehalfspacing

% capitulos del documento
\chapter{Introducción}
\label{chapter:introduccion}


%%% SECTION
\section{Descripción general del problema}
\label{def:def1}

\paragraph{}
El Trabajo de Final de Máster (\textit{TFM}) que aquí se presenta nace de la necesidad por parte de los profesionales sanitarios de poder tener la información más exacta posible sobre las mejores referencias bibliográficas actuales sobre tratamientos a aplicar a un paciente, dado unos síntomas concretos. Actualmente existe infinidad de referencias medicas que los profesionales sanitarios pueden consultar, pero esta información es tan abundante que acaba siendo engorrosa de consultar. Esto hace que, muchas veces sea complicado encontrar la información sobre estas referencias bibliográficas para tratar algunas enfermedades. En el ámbito de la medicina el tiempo perdido puede costar vidas y es un precio demasiado elevado a pagar, tanto a nivel económico como emocional.

\paragraph{}
Los profesionales sanitarios necesitan poder disponer de una plataforma que sea capaz de poder facilitarles estas referencias bibliográficas actuales lo más exactas y personalizadas que sean posible, ajustándose a la búsqueda de la información que poseen de sus pacientes.

\paragraph{}
Bajo esa premisa nace el proyecto HOPE (que significa \textit{Health Operations for Personalized Evidence} en ingles) con el objetivo de ayudar a estos profesionales sanitarios a encontrar estas referencias bibliográficas que necesitan de la manera más rápida y fácil posible. Actualmente existen bases de datos de confianza en donde los profesionales sanitarios y el publico en general puede buscar estas referencias de ensayos y estudios clínicos con información fiable y desarrollados con anterioridad, pero no siempre es fácil o rápido encontrar estos resultados.

\paragraph{}
El proyecto HOPE es un sistema basado en inteligencia artificial para identificar la información de casos clínicos registrados en la Historia Clínica Electrónica, en base a los cuales realiza una búsqueda única por paciente para proporcionar al profesional sanitario recomendaciones de referencias bibliográficas donde constan tratamientos, estudios de investigación e información para ayudar al paciente. Todo en base a registros de fuentes científicas de información. En este proyecto, los profesionales sanitarios de todo el mundo puede consultar en una base de datos estas referencias y ver que otros tratamientos relacionados con los síntomas de sus pacientes han dado resultado. Todo y con eso, el sistema no siempre devuelve las referencias actuales más relevantes por lo que, no siempre la información consultada es útil.

\paragraph{}
Ademas, actualmente los profesionales sanitarios pueden valorar si la información recibida ha sido útil o no respecto a la búsqueda que han realizado, por lo que con esta valoración (\textit{feedback}), se pretende mejorar el sistema actual complementándolo con un modelo clasificador capaz de ayudar al actual a entregar realmente las referencias que son más útiles basándose en la valoración que los profesionales sanitarios dan al sistema.

\paragraph{}
Este Trabajo de final de máster (\textit{TFM}) pretende ayudar a mejorar el proyecto HOPE, mejorando su algoritmo de Inteligencia Artificial para que los resultados se ajusten a las necesidades de información que requieren los profesionales sanitarios, en base a las búsquedas personalizadas que puedan hacer respecto a la información que tienen de sus pacientes. Para hacer esto se realizara un estudio de aproximación a conocer cual es el mejor modelo predictivo que puede ayudar a devolver esa información lo más exacta posible.

\section{Motivación personal}

\paragraph{}
En los duros tiempos en los que estamos viviendo actualmente, tanto económica como emocionalmente, es grato ver como la humanidad es capaz de dejar a un lado sus diferencias y unirse para afrontar problemas comunes. En el caso del proyecto HOPE, lo que más me atrajo fue la oportunidad de poder ayudar a encontrar soluciones a enfermedades que ya están en el último paso (o como médicamente se le describe, en cuidados paliativos).

\paragraph{}
Es evidente que este proyecto no va a dar una solución para curar cualquier enfermedad, pero si puede ayudar a los profesionales sanitarios a poder encontrar posibles soluciones a enfermedades complejas, ayudando a estos profesionales a buscar en la infinidad de documentación medica que existe, el tratamiento que más pueda ayudar a paliar o, quien sabe, curar una enfermedad que ya se daba por incurable. Esto mismo es lo que me motiva y mucho el poder ayudar ante estas situaciones.

\paragraph{}
También me ha motivado muchísimo el conocer a gente profesional que, independientemente del país al que pertenece o la profesión que tiene, se una al proyecto HOPE para participar y ayudar con sus conocimientos a hacer de este mundo, un lugar un poco mejor en el que vivir. Esta experiencia me esta enriqueciendo muchísimo personalmente y espero poder seguir contribuyendo al proyecto cuando este máster acabe.


% bibliografia
\addcontentsline{toc}{chapter}{Bibliografía}
\bibliographystyle{plain}
\bibliography{referencias}

\end{document}